\documentclass[10pt]{article}
\usepackage[T1]{fontenc}
\usepackage[utf8]{inputenc}
\usepackage[french]{babel}
\usepackage{amsmath,amssymb,mathrsfs,mathpazo}
\usepackage{geometry}
\geometry{a4paper,margin=2cm}
\usepackage{graphicx,array,tabularx,tabularray,colortbl,multirow,multicol,booktabs,enumitem,pifont,ulem,fancybox,diagbox,cancel}
\usepackage{pstricks,pstricks-add,pst-plot,pst-tree,pst-node,pst-text,pst-func,pst-math,pst-bspline,pst-3dplot,pst-eucl,pst-solides3d}
\let\pscalculate\undefined
\usepackage{pst-bezier}
\usepackage{tikz,pgfplots,tkz-tab}
\usetikzlibrary{arrows,arrows.meta,patterns,calc,positioning,trees,decorations.markings,shapes.geometric,shapes.misc,shapes.arrows,intersections,decorations.pathmorphing,decorations.pathreplacing,backgrounds,shadows,babel}
\pgfplotsset{compat=1.18}
\usepgfplotslibrary{fillbetween}
\usepackage[np]{numprint}
\usepackage{listings,makecell,pas-tableur,twoopt}
\let\barre\undefined
\usepackage[pstricks]{tablvar}
\definecolor{darkgreen}{RGB}{0,100,0}
\providecommand{\red}{\color{red}}
\providecommand{\blue}{\color{blue}}
\providecommand{\green}{\color{green}}
\providecommand{\black}{\color{black}}
\providecommand{\white}{\color{white}}
\providecommand{\gray}{\color{gray}}
\providecommand{\cyan}{\color{cyan}}
\providecommand{\magenta}{\color{magenta}}
\providecommand{\yellow}{\color{yellow}}
\providecommand{\violet}{\color{violet}}
\providecommand{\orange}{\color{orange}}
\providecommand{\euro}{\texteuro}
\providecommand{\EUR}{\texteuro}
\providecommand{\ieme}{\textsuperscript{e}}
\providecommand{\degres}{\ensuremath{^\circ}}
\providecommand{\degree}{\ensuremath{^\circ}}
\providecommand{\vv}[1]{\overrightarrow{#1}}
\providecommand{\Coord}[1]{#1}
\providecommand{\decosix}{}
\providecommand{\starredbullet}{\ensuremath{\bullet}}
\pagestyle{empty}
\setlength{\parindent}{0pt}
\usepackage{adjustbox}\DeclareUnicodeCharacter{2236}{:}\DeclareUnicodeCharacter{2190}{\ensuremath{\leftarrow}}
\providecommand{\R}{\mathbb{R}}
\providecommand{\N}{\mathbb{N}}
\providecommand{\Z}{\mathbb{Z}}
\providecommand{\Q}{\mathbb{Q}}
\providecommand{\D}{\mathbb{D}}
\providecommand{\C}{\mathbb{C}}
\providecommand{\e}{{\mathrm{e}}}
\providecommand{\vect}[1]{\overrightarrow{#1}}
\providecommand{\vectt}[1]{\overrightarrow{\mathrm{#1}}}
\providecommand{\vvt}[1]{\overrightarrow{\mathrm{#1}}}
\providecommand{\np}[2][]{{#2~#1}}
\providecommand{\vv}[1]{{\overrightarrow{#1}}}
\renewcommand{\d}{{\mathrm{d}}}
\providecommand{\degres}{{^\circ}}
\providecommand{\degre}{{^\circ}}
\providecommand{\blue}{\color{blue}}
\providecommand{\red}{\color{red}}
\providecommand{\enskip}{\quad}
\providecommand{\llbracket}{\lbrack\!\lbrack}
\providecommand{\rrbracket}{\rbrack\!\rbrack}
\providecommand{\euro}{{\text{€}}}
\providecommand{\no}{{\text{n°}}}
\providecommand{\interval}[2]{{[#1;#2]}}
\providecommand{\intervFO}[2]{{[#1;#2[}}
\providecommand{\dotfill}{\ldots}

\providecommand{\ds}{\displaystyle}
\providecommand{\textding}[1]{#1}
\providecommand{\Oij}{\(\mathrm{O};\vec{\imath},\vec{\jmath}\)}
\providecommand{\Oijk}{\(\mathrm{O};\vec{\imath},\vec{\jmath},\vec{k}\)}
\providecommand{\Ouv}{\(\mathrm{O};\vec{u},\vec{v}\)}
\providecommand{\intervFF}[2]{\ensuremath{[#1;#2]}}
\renewcommand{\euro}{\ifmmode\text{\texteuro}\else\texteuro\fi}
\DeclareUnicodeCharacter{00B0}{\ifmmode^{\circ}\else\textdegree\fi}
\DeclareUnicodeCharacter{20AC}{\ifmmode\text{\texteuro}\else\texteuro\fi}
% Énoncés : Éducation nationale. Transcription et corrigés : APMEP, Denis Vergès et contributeurs. Découpage et classement : CoopMaths, Sébastien Lozano, Éric Elter et contributeurs. Adaptation de la présentation : Maths Explorer.
% CC BY-SA : https://coopmaths.fr/www/about/
% https://coopmaths.fr/alea/static/bac/2025/tex/bac_2025_06_sujet1_etrangers_3.tex
\begin{document}
\noindent\textit{APMEP / CoopMaths. CC BY-SA. Présentation adaptée par Maths Explorer.}\par\medskip
\subsection*{Partie A}

On considère la fonction $f$ définie sur l'intervalle $] - 1~;~+\infty[$ par 

\[f(x) = 4\ln(x + 1) - \dfrac{x^2}{25}\]

On admet que la fonction $f$ est dérivable sur l'intervalle $] - 1~;~+\infty[$.

\medskip

\begin{enumerate}
\item Déterminer la limite de la fonction $f$ en $-1$.
\item Montrer que, pour tout $x$ appartenant à l'intervalle $]- 1~;~+\infty[$ , on a :
 \[f'(x) = \dfrac{100- 2x - 2x^2}{25(x+1)}\]

\item Étudier les variations de la fonction $f$ sur l'intervalle $] - 1~;~+\infty[$ puis en déduire que la fonction $f$ est strictement croissante sur l'intervalle [2~;~6,5].
\item On considère $h$ la fonction définie sur l'intervalle [2~;~6,5] par $h(x) = f(x) - x$.

On donne ci-dessous le tableau de variations de la fonction $h$ :

\begin{center}
\psset{unit=1cm,arrowsize=2pt 3}
\begin{pspicture}(11,2.5)
\psframe(11,2.5)\psline(0,2)(11,2) \psline(1,0)(1,2.5)
\uput[u](0.5,1.9){$x$}\uput[u](1.2,1.9){$2$}\uput[u](6,1.9){$m\approx 2,364$}\uput[u](10.75,1.9){$6,5$}
\rput(0.5,1){$h(x)$}\uput[u](1.4,0){$h(2)$}\uput[d](6,2){$M \approx 2,265$}\uput[u](10.5,0){$h(6,5)$}
\psline{->}(1.8,0.4)(4.9,1.6)\psline{->}(6.9,1.6)(10,0.4)
\end{pspicture}
\end{center}
Montrer que l’équation $h(x) = 0$ admet une unique solution $\alpha \in [2~;~6,5]$.

\item On considère le script suivant, écrit en langage Python:

\begin{center}
\begin{tabular}{|l|}\hline
\quad \textbf{from} math \textbf{import *}\\
~\\
\quad \textbf{def} f(x):\\
\quad \qquad\textbf{return} 4*log(1+x)-(x**2)/25\\
~\\
\quad \textbf{def} bornes(n):\\
\quad \qquad p = 1/10**n\\
\quad \qquad x = 6\\
\quad \qquad \textbf{while} f(x)-x > 0:\\
\quad \qquad\qquad x = x + p\\
\quad \qquad\textbf{return} (x-p, x)\\\hline
\end{tabular}
\end{center}

On rappelle qu'en langage Python:


\begin{itemize}[label=$\bullet~$]
\item la commande $\text{log} (x)$ renvoie la valeur $\ln x$ ; 
\item la commande c$**$d renvoie la valeur de $\text{c}^{\text{d}}$.
\end{itemize}

	\begin{enumerate}
		\item Donner les valeurs renvoyées par la commande bornes(2).
		
On donnera les valeurs arrondies au centième.
		\item Interpréter ces valeurs dans le contexte de l'exercice.
	\end{enumerate}
\end{enumerate}

\subsection*{Partie B}

\emph{Dans cette partie, on pourra utiliser les résultats obtenus dans la partie A.}

On considère la suite $(u_n)$ définie par $u_0 = 2$, et, pour tout entier naturel $n,\: u_{n+1} = f\left(u_n\right)$

\begin{enumerate}
\item Montrer par récurrence que pour tout $n$ entier naturel,

\[2 \leqslant u_n \leqslant u_{n+1} < 6,5.\]
\item En déduire que la suite $(u_n)$ converge vers une limite $\ell$.
\item On rappelle que le réel $\alpha$, défini dans la partie A, est la solution de l'équation $h(x)~=~0$ sur l'intervalle [2~;~6,5].

Justifier que $\ell = \alpha$.
\end{enumerate}

\bigskip

\end{document}
